% Proof of Lemma \ref{lem:isotropy-moments}
\begin{proof}
\textbf{(1):} Take $B = e_i e_j^\top$ in the isotropy condition. The left side is $\E[X e_i e_j^\top X^\top] = \E[X_{\cdot j} X_{\cdot j}^\top]$, which has $(k, l)$ entry $\E[X_{kj} X_{lj}]$. The right side is $\frac{\sigma^2}{N} \cdot 1 \cdot I = \frac{\sigma^2}{N} I$. For $k = l = i$, we get $\E[X_{ij}^2] = \sigma^2/N$. For off-diagonal, the isotropy condition at different choices of $B$ yields the centered condition.

\medskip
\noindent\textbf{(2):} This follows from expanding the isotropy condition with $B = e_i e_k^\top$ and comparing entries.

\medskip
\noindent\textbf{(3):} We have
\begin{align}
\E[\Tr(XBX^\top C)]
&= \E\Bigl[\sum_{ijkl} X_{ij} B_{jk} X_{il} C_{li}\Bigr] \\
&= \sum_{ijkl} B_{jk} C_{li} \E[X_{ij} X_{il}] \\
&= \sum_{ijkl} B_{jk} C_{li} \frac{\sigma^2}{N} \delta_{jl} \\
&= \frac{\sigma^2}{N} \sum_{ij} B_{jj} C_{jj} = \frac{\sigma^2}{N} \Tr(B) \Tr(C).
\end{align}
The $O(1)$ term arises from boundary effects when $B$ or $C$ have rank $O(1)$.
\end{proof}
